% ===== Creation and Co-Creation: the unity of aesthetics and production =====
\section{Creation and Co-Creation: The Unity of Aesthetics and Production}
\label{sec:creation}

\noindent
The previous section argued that the construction of the field is an aesthetic and not an engineering problem, and in doing so it used a word, ``construction'', that must now be examined and replaced, for it carries an engineering residue the argument has been straining against. To build, to construct, suggests a making from a plan: the prior design, the executed assembly, the \emph{poiesis} that brings forth according to a blueprint. But what the field's maker does is not this. It is closer to what an artist does than to what a builder does, and the proper word for it is \emph{creation} rather than construction, the bringing-into-being of something that did not exist and was not specified in advance, the opening of a world rather than the assembly of a structure. And once the word is corrected, a concept enters that the paper, having reduced travel to a field and the field to an aesthetic matter, can no longer avoid: for the central act of the aesthetic is creation, and a philosophy that has made the aesthetic its centre must give an account of what creation is. This section gives that account, and the account has a consequence that reorganises the whole paper: under the relational ontology, creation is \emph{co-creation}, and co-creation is the same activity the paper has called, in its political economy, relational production. The aesthetic concept (creation) and the political-economic concept (production) are revealed as one, and their unity is the keystone that joins the paper's two great lines.

% ============================================================
\subsection{From construction to creation}
\label{sub:cr-fromconstruction}

The reasons for preferring ``creation'' to ``construction'' are not merely terminological; they track a real difference in the kind of act involved. Construction, in its proper sense, is the realisation of a design: one knows what is to be built, and builds it; the end precedes the making as a plan precedes its execution, and the maker's skill is the faithful translation of the plan into the thing. This is exactly the engineering model the previous section rejected (\xref{sub:constr-paradox}): if the field were constructed, it would be specified in advance and assembled to specification, and the contingency, spontaneity, and openness that the field requires would be designed out of it. Creation is the other thing. To create is to bring forth what was not specified in advance, to open a space in which something genuinely new can appear, something the creator did not and could not fully foresee, since if it were foreseen in full it would be executed rather than created. The artist does not execute the poem from a complete prior specification of it; the poem comes to be in the making, its final form discovered rather than implemented, the creation being precisely the bringing-forth of what exceeds any prior plan. This is why the construction of the field is better called its creation: what the field's maker does is open a space in which a genuinely new relational reality, the lived ``we'' of the journey, with all its unforeseen contingency, can come to be, and this opening of space for the unforeseen is creation, not construction.

The correction matters because it brings the field's making under the concept that the aesthetic tradition has most deeply theorised. The making of beauty is creation: the work of art is created, not constructed; and if the construction of the field is an aesthetic matter, then it is a matter of creation, and the field's maker is, in the relevant sense, an artist, one who creates a space for beauty rather than one who builds a structure to specification. Everything the paper has said about the field's making being aesthetic (\xref{sec:construction}) is thus gathered into the single claim that the field is \emph{created}; and the philosophy of beauty, which the next sections develop, is therefore also, at its root, a philosophy of creation, for the beautiful is the created, and the question of beauty is inseparable from the question of how the beautiful comes to be.

% ============================================================
\subsection{Under the relational ontology, creation is co-creation}
\label{sub:cr-cocreation}

The concept of creation enters the paper already inflected by the tradition that theorised it, and that tradition's default must be confronted, for it is at odds with everything the paper has established. The Western theory of creation is, characteristically, a theory of the \emph{single} creator. The paradigm is the solitary artist: the genius who, in Kant's account, is the one through whom nature gives the rule to art, the singular individual whose originality is the source of the new \citep{kant1987}; the Romantic creative imagination, the lone author from whom the work issues. Creation, on this paradigm, is the act of one subject who, facing the material, brings forth from within (or from the muse, or the unconscious) an object that did not exist. It is solitary: the creator confronts the world and produces the work, and the work is the creator's.

But under the relational ontology this paradigm is, if not simply false, then derivative, a special and impoverished case of something whose full form is relational. For if the subject is itself relational, constituted in and by relation (\xref{sub:ap-relationalsubject}), then creation, as an act of the subject, is no more originally solitary than judgement or perception were. Just as the paper showed that aesthetic judgement is originally relational, completed in the shared ``look'' and merely truncated in the solitary aesthete (\xref{sub:ap-relationalsubject}), so creation is originally \emph{co-creation}, the joint bringing-forth by a relational subject, and merely truncated in the solitary creator. The solitary artist is not the paradigm of creation from which joint creation is a derivative complication; the solitary artist is the \emph{deficient} case, creation arrested at one pole of a relation, and the full form of creation is the co-creation in which two (or more) bring forth together what neither could bring forth alone. This is not a romantic exaggeration of collaboration but a consequence of the relational ontology: if the creating subject is relational, then creation is the act of a relational subject, which is to say it is co-creation, and the lone genius is the limiting case in which the relational structure of creation has been narrowed to a single pole.

The field of travel makes this vivid, for the field is unmistakably co-created (\xref{sub:field-cocreation}): it is not created by one traveller for the other, as an artist makes a work for an audience, but brought forth by the two together, each contributing to the opening of a space that becomes the ``we'''s and neither's alone. The journey co-imagined before departure, the field co-opened in the travelling, the experience co-generated in the living, all are co-creation, the joint bringing-forth of a relational reality by the relational subject whose reality it is. And what is vivid in travel is, the paper claims, the truth of creation as such: creation is co-creation, the bringing-forth of the new by a relational subject, of which the solitary artist's making is the narrowed, single-poled case. The deepest creation, as the deepest aesthetic judgement, the deepest love, is the one conducted together.

% ============================================================
\subsection{Co-creation is relational production: the unity of the two lines}
\label{sub:cr-unity}

The recognition that creation is co-creation yields the section's central result, which is one of the most important in the paper: that co-creation, the aesthetic concept, and relational production, the political-economic concept, are \emph{the same activity under two names}. The paper has run, throughout, on two parallel lines. One is aesthetic: the creation of the field, the perception of beauty, the making of the beautiful. The other is political-economic: the production of experience, the reproduction of relational wealth, the generation of value by the ``we'' (\xref{sec:reproduction}, \xref{sec:relational-production}). These have seemed distinct vocabularies for distinct aspects. They are not. Co-creation and relational production name one activity, \emph{the bringing-forth of the new by a relational subject}, seen from two sides.

Seen from the aesthetic side, the activity is \emph{co-creation}: two relational subjects bringing forth together something that did not exist, a field, an experience, a meaning, a beauty. Seen from the political-economic side, the same activity is \emph{relational production}: a relational subject generating and reproducing the value proper to the relation. The aesthetic emphasis falls on the novelty, the bringing-forth of what did not exist, the creative opening; the political-economic emphasis falls on the value generated, its reproduction and circulation, its justice. But the activity is one: in both, a relational subject brings forth, together, what neither member could bring forth alone, and what it brings forth is at once a new thing (the aesthetic accent) and a value (the political-economic accent). This is why the paper's two lines have converged again and again, why the reproduction of experience turned out to be the evolution of aesthetic judgement (\xref{sec:evolution}), why the good cycle of value and the good cycle of beauty were the same cycle: because the aesthetic and the political-economic were never two activities but one, co-creation and relational production being two names for the bringing-forth of the new and valuable by a relational subject.

This unity is anchored in the deepest stratum of the materialist tradition, which the paper here makes explicit. For Marx, production was never merely economic: in its root sense it is the human being's objectification of itself in a world it creates, the activity by which the human, as a ``species-being,'' produces an objective world and in producing it produces itself, creates its own life and its own nature through the creative transformation of the material world \citep{marx1988}. Production, in this foundational sense, \emph{is} creation, the creative objectification of the human in a world of its own making, and the later, narrower, economic sense of production is a specification of this creative root. So the convergence of creation and production is not the forcing-together of two alien concepts; it is the recovery of their original unity, which the materialist tradition had grasped at its foundation: production is the materialist and social name for creation, and creation is the aesthetic name for production, and both name the bringing-forth by which a subject objectifies itself in a world. Under the relational ontology, where the subject is relational, this bringing-forth is co-creation \emph{is} relational production: the relational subject objectifies itself in the world it brings forth together, creating, in creating its shared world, itself as a ``we.'' The two lines of the paper are one line, and creation is the name of their unity.

A nuance preserves the accents while affirming the unity. Creation and production are the same activity, but the words carry different emphases worth keeping: ``creation'' foregrounds the bringing-forth of the genuinely new, the unforeseen, the opening of what did not exist; ``production'' foregrounds the generation and reproduction of value, its circulation and justice. The same act of a relational subject is creation under the first emphasis and production under the second, and the paper will use whichever word foregrounds the aspect in view, but the reader should hold that they name one activity, the relational subject's joint bringing-forth, which is at once creative (it makes the new) and productive (it generates value), aesthetic (it makes beauty) and political-economic (it makes wealth). It is on this unified concept, co-creation as relational production, the bringing-forth of the new and valuable by a relational subject, that the ethics of creation now builds.

% ============================================================
\subsection{The ethics of creation: co-created works and the question of destruction}
\label{sub:cr-ethics}

If creation is co-creation and co-creation is relational production, then the ethics of creation falls under the ethics of relational production already established (\xref{sec:relational-production}), but it adds something the earlier treatment did not make explicit: that relational production has a \emph{negative} as well as a positive moment. The positive moment is creation, the bringing-forth of the new; the negative moment is \emph{destruction}, the unmaking of what was brought forth. The ethics of relational production must govern both, and the question of destruction, of when, and by whom, a created thing may rightly be destroyed, is the sharp test of the whole framework, for it is where the principle of subject-preservation is most concretely at stake.

The question has a familiar answer under the solitary paradigm of creation, and the answer is wrong here for an instructive reason. Under the solitary paradigm, the creator owns the creation and may destroy it: Kafka may will his manuscripts burned; the painter may slash the canvas; the maker, having made, may unmake, for the work is the maker's and the maker's alone. This treats creation on the model of production-as-possession: I made it, therefore I own it, therefore I may destroy it. And for genuinely solitary creation the intuition has some force. But the works of relational production are not solitary creations; they are \emph{co-created}, and for a co-created thing the destruction-right dissolves. Consider the sharp and ordinary case: in a relationship, in a fit of anger, one partner deletes the shared photographs, the images of a shared life, co-created, the very wealth of the ``we'' (\xref{sec:reproduction}). The act is a wrong, and the framework says precisely why. The photographs are the co-created wealth of the relational subject rather than one partner's possession; to destroy them unilaterally is to \emph{annihilate rather than to dispose of one's own creation, single-handedly, a creation of the ``we.''} It violates every criterion of relational production at once: the ontological (it is unilateral, where what belongs to the ``we'' may be unmade only by the ``we'', destruction, the negative of reproduction, must like reproduction be joint); the teleological (it forecloses the future in which that shared experience might have been reproduced, deepened, transfigured, slamming shut a path of the relation's continued generativity); the political-economic and the power criterion together (in anger, the destruction is wielded as a weapon, the shared wealth converted into an instrument of harm and domination). And over all of these stands the keystone (\xref{sub:rp-subjectpreservation}): the unilateral destruction of a co-created thing annihilates the other's part in it, the other's memory, investment, presence in the shared creation, and so purchases one partner's momentary unity-with-their-rage at the cost of annihilating the other's share in the ``we.'' There is, the framework concludes, \emph{no unilateral right to destroy a co-created thing}; destruction, like creation, must be relational, conducted by the ``we'' whose creation it is, on pain of violating the principle of subject-preservation.

Three tensions must be addressed, for the matter is not as simple as a flat prohibition. The first: what if the partners contributed unequally, if the photographs were mostly one partner's making, cherished by one and ignored by the other? This grants no unilateral destruction-right, and the reason discloses something deep about relational production: the value of a co-created thing does not belong to the partners severally, in proportion to their contributions, but to the ``we'' as such, for to parcel out the relation's wealth by individual contribution is already to dissolve the relational subject into two calculating individuals, which is itself the corruption the theory forbids (\xref{sub:rp-category}). Unequal contribution does not convert co-created wealth into private property; it remains the ``we'''s, and unmakeable only by the ``we.'' The second tension: the Kafka case, of genuinely solitary creation. Here the framework gives a \emph{graded} rather than an absolute answer. Even ``one's own'' creation is rarely purely solitary under the relational ontology, for it stands in relation to others, readers, inheritors, the culture it enters (\xref{sub:rp-social}), and once it has entered those relations it is no longer purely private to destroy. But this relational claim on a solitary creation is \emph{weaker} than the claim a co-created thing carries: the destruction of a purely personal creation is a tension between a real (if not absolute) personal right of disposal and a responsibility to the relations the creation has entered, whereas the destruction of a co-created thing is a clear prohibition, no unilateral right existing at all. The framework yields a gradient: the more a creation is co-created, or has entered into relations, the weaker any unilateral right to destroy it; the more purely solitary and unentered, the stronger the personal right of disposal, but since purely solitary, unentered creation is a limiting case that the relational ontology makes nearly empty, an absolute unilateral destruction-right is nearly never established. The third tension: destruction is not always a wrong, for destruction can itself be creative, the artist who paints over a failed canvas to free the next, the Daoist diminishing (\zh{为道日损}), the clearing-away that makes room for new growth, even in a relationship the letting-go of what must be released to begin again. The framework distinguishes good from bad destruction by the same criteria: good destruction preserves the possibility of continued generativity (it unmakes in order that something new may be made), is joint where the destroyed thing is co-created, and is neither exploitative nor wielded as domination; bad destruction forecloses generativity (it ends rather than renews), is unilateral over a co-created thing, and is wielded as a weapon. Destruction as such is not the wrong; the wrong is the \emph{unilateral, generativity-foreclosing, weaponised destruction of a co-created thing}. The deletion in anger is wrong on every count; the mutual decision to let some painful relic go, freeing the relation for new creation, may be entirely good. The criteria discriminate, as throughout, not by the act's surface but by whether it preserves or annihilates the relational subject and its generativity.

The ethics of creation is therefore symmetrical, governing a positive moment (co-creation, the bringing-forth of the new) and a negative moment (destruction, the unmaking of what was brought forth), and the principle of subject-preservation reigns over both: as the relation may not create its unity by annihilating a subject, so it may not destroy a co-creation by the unilateral act of one, for the unmaking of what the ``we'' made, like its making, belongs to the ``we.'' Destruction, like creation, is relational; and the lone hand that unmakes the shared work has, in the unmaking, annihilated the other's share in it.

% ============================================================
\subsection{Creation, the public sphere, and the limit of the theory}
\label{sub:cr-public}

There is a further case that presses beyond the relational frame, and it must be met not by extending the theory to cover it but by marking, honestly, where the theory ends. A created work does not only enter into particular relations; it may enter the \emph{public sphere}, may become public culture, and once it has, it belongs no longer to its author alone, nor, the paper must concede, to any determinate ``we'' at all. Kafka's manuscripts, once they have become \emph{The Castle}, once they are world literature, are no longer his to burn without remainder, for a public has a claim on them; but the public that has this claim is not a relation, not a particular ``we'' of nameable others, but an anonymous, indefinite, open-ended domain, including readers not yet born, a ``public'' that is precisely not a relational subject. This case raises, in an acute form, a question about the reach of the theory of relational production, and the question must be faced directly, for there is a temptation to which the theory must not yield.

The temptation is to interpret the public sphere as a relational domain, to say that the public is simply a larger ``we,'' that public culture is a vast co-creation, and so to bring the public sphere under the theory of relational production whole. This would be a mistake, and naming it is important, for it is the characteristic over-reach of a successful theory: the imperialism by which a concept that illuminates one domain claims to illuminate all. The public sphere is \emph{not} a magnified relation, and to interpret it as one would diminish rather than extend the analysis, for it would erase exactly the features that constitute the public sphere as public. The public sphere is anonymous where relation is particular: it is addressed to indefinite, unnameable others, including the unborn, not to the specific other of a relation \citep{habermas1989}. It is non-reciprocal where relation is mutual: one may owe something to posterity, but posterity cannot answer, and the reciprocal structure of the relational subject (\xref{sub:ap-relationalsubject}), the mutual ``look,'' has no purchase on a public that cannot look back. It is institutional and durable where relation is intimate: the public realm has its own logic of appearance, permanence, and the common world (\citealp{arendt1958}), its own institutions of law, heritage, and publicity, which are not the logic of intimate relation enlarged. And it is the contested field of political economy (\xref{sec:pe-beauty}), traversed by power and capital in ways that are not reducible to the good and bad cycles of relational subjects. To treat the public sphere as a big relation would abolish all of this, would shrink the public to the intimate, the anonymous to the particular, the non-reciprocal to the mutual, and so would not enlarge the theory's explanatory reach but contract its object, dissolving the public sphere's own reality into a relational image that misses what is essential to it.

The right course is therefore to mark a \emph{limit}, and in marking it to say precisely what the theory of relational production does and does not explain about a work that has entered the public. What it explains is the work's persisting \emph{relational-productive dimension}: even in the public sphere, the work remains a created thing, a product of co-creation or solitary creation, bearing relational value, and its destruction still touches the question of subject-preservation insofar as it remains anyone's creation. This dimension the theory illuminates, and it does not vanish when the work goes public; a co-created work that enters public culture still carries the prohibition on its unilateral destruction by one of its co-creators, now reinforced by the public's claim. But the public sphere \emph{as} public, its anonymity, its non-reciprocity, its institutionality, its distinctive politics, its mode of belonging to the unborn, is not explained by the theory of relational production and must not be pretended to be. It belongs to other theories: to the theory of the public sphere and the space of appearance \citep{habermas1989,arendt1958}, to the theory of cultural heritage and the commons, to a theory of intergenerational justice, to the political economy the paper has only touched (\xref{sec:pe-beauty}). The relation between the relational-productive dimension of a work and its public dimension, the way a created thing is at once the product of a relational subject and a member of a public world, lying in the overlap of two domains governed by different logics, is a real and important problem, but it is one the theory of relational production opens rather than closes.

This honest marking of the limit is a condition of its rigour rather than a weakness of the theory, and it is worth stating as a principle of method. A theory that knew no limit, that claimed to explain the public sphere and the intimate relation and all between by the single logic of relational production, would for that very reason explain none of them well, for it would have dissolved the differences that make each what it is. The power of the theory of relational production in its own domain, the domain of the particular relation, the ``we'' of nameable others, the co-creation and reproduction of shared value, depends on its not being stretched to a domain with a different logic. The theory illuminates the relational, and the relational dimension of what enters the public; it reaches its limit where the public sphere's own, non-relational features begin, and it points, there, beyond itself. The full ethics of a work's life in the public sphere, the overlap and the difference of the relational-productive and the public, the destruction and preservation of public creations, the claims of the anonymous and the unborn, is accordingly named here not as a result but as a directed opening (\xref{sub:clo-fissures}), a problem the theory of relational production frames and hands on to a theory of the public it does not itself contain. To know this limit is to know what relational production is, and what it is not, which is the beginning of using it well.
